\documentclass[11pt]{article}

% ---------- Expert 5: Document Architect — preamble ----------
\usepackage[utf8]{inputenc}
\usepackage[T1]{fontenc}
\usepackage[margin=1in]{geometry}
\usepackage{amsmath,amssymb}
\usepackage{booktabs}
\usepackage{natbib}
\usepackage{enumitem}
\usepackage[T1]{fontenc}\usepackage[protrusion=true,expansion=false]{microtype}
\usepackage{xcolor}
\usepackage[colorlinks=true,linkcolor=blue,citecolor=blue,urlcolor=blue]{hyperref}
\usepackage{titlesec}

\titleformat{\section}{\normalfont\Large\bfseries}{\thesection}{1em}{}
\titleformat{\subsection}{\normalfont\large\bfseries}{\thesubsection}{1em}{}

\newcommand{\rhoDE}{\rho_{\mathrm{DE}}}
\newcommand{\Om}{\Omega_m}

\title{\bfseries A Methodology for AI-Assisted Computational Cosmology:\\
Process, Comparison, and Adversarial Validation\\[0.4em]
\large A Retrospective on the Development of a Dark-Energy\\
Feature-Detectability Study}

\author{%
Documented by a five-role authoring panel:\\
Research Methodologist, AI/Cognitive Scientist, Ethics \& Reproducibility Auditor,\\
Red-Teaming Analyst, and Document Architect\thanks{This paper documents the
\emph{process} of a research collaboration between a human domain researcher and
an AI assistant. It is a methodology retrospective, not the scientific paper
itself; the underlying physics results are reported separately.}}

\date{\today}

\begin{document}
\maketitle

\begin{abstract}
We document the methodology of a human--AI collaborative research project in
observational cosmology, whose scientific output was a study of whether
expansion-history probes can detect and characterize a sudden transition
(``step'') in the dark-energy density. Rather than reporting the physics, this
paper reports the \emph{process}: the chronological path from an informal
brainstorm to a validated manuscript, the specific breakthroughs and reversals
that reshaped the work, and a candid comparison of the AI-assisted workflow
against conventional research practice. We give particular attention to failure
modes specific to AI-assisted science---confirmation-bias feedback loops, model
hallucination, and alignment or ``sycophancy'' drift---and to the concrete
red-teaming protocols used to guard against them. We argue that the central
methodological lesson is not that AI accelerates research, but that it changes
\emph{where} the researcher's vigilance must be spent: from execution toward
continuous, adversarial verification. We close with a limitations analysis,
a broader-impacts discussion, and a reproducibility statement.
\end{abstract}

\tableofcontents
\bigskip

% ============================================================
\section{Introduction and Scope}
\label{sec:intro}

This document is a methodology retrospective. The scientific project it describes
asked a well-posed question in observational cosmology: if the dark-energy density
$\rhoDE(z)$ underwent a sudden change at some redshift---rather than evolving
smoothly---which cosmological probes could detect it, and what could they measure
about it? That question, its answer, and the supporting computations are reported
in the companion physics manuscript. Here we instead examine \emph{how the work
was done}, because the collaboration was unusual in one respect: the bulk of the
implementation, literature synthesis, debugging, and drafting was carried out by
an AI assistant working interactively with a human researcher, over an extended
sequence of sessions.

We believe the process is worth documenting for three reasons. First, the project
exhibited a clear, instructive arc in which several early conclusions were
\emph{overturned} by later, more careful work---a pattern that reveals how
AI-assisted research can go wrong and how it can be made to self-correct. Second,
the division of cognitive labor between human and machine differed markedly from
both traditional solo research and traditional collaboration, in ways that suggest
transferable practices. Third, the failure modes we encountered are precisely those
the AI-safety and research-integrity communities have flagged as systemic risks,
and we can report concretely on which materialized and which were successfully
mitigated.

The intended audience is deliberately broad: practitioners in observational
cosmology who may wish to adopt or critique the workflow; AI researchers and
enthusiasts interested in a concrete, non-trivial scientific application; and
research-integrity specialists concerned with the epistemics of machine-assisted
discovery.

% ============================================================
\section{Expert 1 --- The Research Methodology and Its History}
\label{sec:history}

\subsection{Origin: an informal analogy, disciplined into a question}

The project did not begin with a hypothesis. It began with a loose analogy---the
question of whether dark energy could ``discharge'' or ``crack'' in a manner
loosely reminiscent of a sudden physical release, generated in an exploratory,
free-associative brainstorm. The first methodological act was therefore
\emph{triage}: separating the physically defensible kernel of the idea from the
parts that were metaphor. Dark energy cannot radiate a detectable gravitational-wave
signature from such an event (any such signal redshifts far below the band of
foreseeable detectors), so the only observational handle is the background
expansion history $H(z)$. This reduction---from an evocative image to a precise,
testable statement about a step in $\rhoDE(z)$---set the entire subsequent
program.

The lesson, recorded early, is that the value of the initial analogy was not its
correctness but its \emph{generativity}: it pointed at a question
(``can we detect an abrupt dark-energy feature?'') that turned out to be
under-explored, even though the specific mechanism that inspired it was not viable.

\subsection{The methodological spine}

The work settled quickly on a fixed methodological template that was reused for
every probe and every variant:
\begin{enumerate}[leftmargin=1.5em]
  \item define a parametric feature model (a $\tanh$ step in $\rhoDE(z)$ with
        amplitude $A$, location $z_t$, width $\Delta z$);
  \item embed it in the Friedmann equation and forward-model each probe's
        observables;
  \item perform a nested likelihood-ratio test on a dense parameter grid;
  \item \textbf{validate the null distribution} against the analytic
        $\chi^2_1$ expectation before trusting any detection rate;
  \item report Monte-Carlo detection probabilities from the real covariance of
        each probe.
\end{enumerate}
Step~4 is the load-bearing one. It is the discipline that repeatedly caught errors,
and it is the single most transferable element of the methodology.

\subsection{The first breakthrough: a bug that masqueraded as a result}

An early implementation used a multi-start numerical optimizer to fit the nested
models. It produced \emph{negative} likelihood-ratio statistics---a restricted
model apparently fitting better than the model that contained it, which is
impossible. Rather than patch the number, the team treated the impossibility as
diagnostic. The root cause was that the optimizer failed to enforce nesting: the
step model has an ill-defined location and width when its amplitude is zero (a
Chernoff-type boundary problem), so the null was not $\chi^2_3$ and the optimizer
wandered. The fix---replacing the optimizer with a dense grid that includes
$A=0$ exactly, guaranteeing nesting---made the null match $\chi^2_1$ to two
decimal places. \textbf{The methodological principle crystallized here:} in
nested model comparison, a negative $\Delta\chi^2$ is not noise to be clipped but
a signal that the inference machinery is broken.

\subsection{The pivotal reversal: synthetic optimism meets real data}

The most consequential episode was a reversal. Initial detection rates were
computed with a \emph{synthetic} supernova survey carrying a realistic error
budget. These suggested supernovae could detect a large dark-energy step a
substantial fraction of the time. When the same analysis was rerun against the
\emph{real} Pantheon+ compilation with its full published stat+sys covariance,
the detection rate collapsed by roughly an order of magnitude. The synthetic
survey had been optimistic in exactly the way synthetic surveys usually are: it
lacked the correlated systematics and the parameter degeneracies of the real
instrument.

This forced a deeper question---\emph{why}---and produced the project's central
scientific insight, which in turn only emerged because the reversal demanded an
explanation. A single integrated distance (a supernova magnitude) is nearly
degenerate with the matter density $\Om$: both reshape $H(z)$ similarly, so a
floating $\Om$ absorbs the step. Anisotropic BAO escapes this because it supplies
\emph{two} geometric observables at each redshift, one of which (the radial
distance, $\propto 1/H$) responds to the step directly. The mechanism was then
proven by a controlled numerical experiment: stripping BAO to a single distance
collapsed its sensitivity to supernova-like blindness.

\subsection{Consolidation and extension}

With the mechanism understood, the same validated apparatus was turned to a
sequence of further questions, each a small, self-contained study: what amplitudes
the real data already \emph{exclude}; where in redshift each probe is most
sensitive; whether the transition \emph{width} is recoverable (it is not, at
current precision); and whether the framework applies to a physically motivated
model rather than a phenomenological step. The last of these led to an extension
into horizon-entropy (holographic) dark energy, where the same
degeneracy-breaking logic reappeared: a smooth deformation could be constrained
only when an early-universe anchor (the CMB) was added. Finally, the accumulated
results were unified under a single, standard criterion from the Fisher formalism
(feature recoverability as a subspace-overlap condition), explicitly framed as an
\emph{application} of existing theory rather than a novel invention.

\subsection{Chronological summary}

\begin{table}[h]
\centering
\caption{Condensed timeline of the research process, emphasizing the points at
which the direction of the work changed.}
\label{tab:timeline}
\begin{tabular}{@{}p{0.24\textwidth}p{0.68\textwidth}@{}}
\toprule
Phase & Character and turning points \\
\midrule
Brainstorm & Informal analogy; triaged to a precise question about a step in $\rhoDE(z)$. \\
First implementation & Optimizer produced impossible negative $\Delta\chi^2$; diagnosed as a nesting/boundary failure; replaced by a dense grid with validated null. \\
Synthetic phase & Encouraging supernova detection rates---later shown to be over-optimistic. \\
\textbf{Reversal} & Real Pantheon+ covariance collapsed supernova sensitivity $\sim$10$\times$; demanded a mechanistic explanation. \\
Mechanism & Two-observable structure of BAO breaks the step--$\Om$ degeneracy; proven by a controlled single-observable experiment. \\
Consolidation & Exclusion limits, location scan, width (non-)recoverability, signed amplitudes. \\
Physical extension & Holographic (Barrow/Tsallis/R\'enyi) dark energy; CMB distance priors; multi-model comparison. \\
Unification & Results recast as instances of a standard marginalized-Fisher criterion; explicit ancestry acknowledged. \\
Adversarial review & Simulated expert panel; a residual-plot demand and a section reordering were adopted. \\
\bottomrule
\end{tabular}
\end{table}

% ============================================================
\section{Expert 2 --- Comparative Analysis: AI-Assisted vs.\ Conventional Practice}
\label{sec:comparison}

\subsection{What was structurally different}

The collaboration inverted the usual time-allocation of computational research. In
conventional practice, a researcher spends the majority of effort on
\emph{execution}---writing code, wrangling data formats, debugging, formatting
figures---and a minority on \emph{judgment}: deciding what to compute, whether to
believe a result, and what it means. The AI-assisted workflow compressed execution
to near-zero marginal cost and thereby shifted almost the entire cognitive burden
onto judgment. Implementing a new probe, a new estimator, or a new figure became a
matter of minutes, which meant the rate-limiting resource was no longer typing but
\emph{deciding what deserved to be typed and whether to trust the output}.

This has a subtle consequence that we want to state plainly for practitioners: the
cheapness of execution is a double-edged instrument. It enables an unusually
thorough exploration---many small, self-contained sub-studies that a
time-constrained human might never attempt---but it also makes it trivially easy to
generate plausible-looking results faster than they can be checked. The
productivity gain is real only if verification scales with generation. When it does
not, the workflow manufactures unearned confidence.

\subsection{Actionable practices human researchers can adopt}

Several elements of this workflow are separable from the AI and can be adopted by
any researcher:

\begin{enumerate}[leftmargin=1.5em]
  \item \textbf{Validate the null before trusting the signal.} Every detection
        statistic in this project was checked against its analytic null
        distribution before any rate was reported. This single habit caught the
        most serious methodological error. It requires no AI and is undervalued in
        routine practice.
  \item \textbf{Treat impossibilities as diagnostics.} A negative $\Delta\chi^2$,
        a $\chi^2/\text{dof}$ far from unity, a parameter railing against a
        boundary---these are not to be smoothed over. Each such anomaly in this
        project pointed at a specific bug.
  \item \textbf{Prefer the controlled experiment to the plausible narrative.} The
        two-observable mechanism was not asserted; it was \emph{proven} by
        artificially removing one observable and watching sensitivity collapse.
        When an explanation is cheap to test directly, test it directly.
  \item \textbf{Re-run on real data early.} The synthetic-to-real reversal would
        have been far more costly had it come late. The cost of premature reliance
        on idealized inputs is a general hazard, not an AI-specific one.
  \item \textbf{Separate generativity from correctness in early ideas.} The seeding
        analogy was wrong in its specifics but valuable in its direction. Human
        researchers often discard generative-but-flawed ideas too early or commit
        to them too hard; the discipline is to mine the question and discard the
        mechanism.
  \item \textbf{Write the limitations while the work is fresh.} Caveats were
        recorded continuously rather than reconstructed at the end, which kept the
        claims calibrated to what the computations actually supported.
\end{enumerate}

\subsection{Where the AI did not substitute for the human}

It is equally important to record what the automation did \emph{not} do. It did not
choose the questions; the decision to pursue the $\rhoDE$-step framing, to re-run on
real data, and to stop extending rather than over-claim were human judgments. It did
not, on its own, resist the temptation to over-interpret; that restraint had to be
imposed and repeatedly reasserted. And it did not supply ground truth: every number
it produced was, in principle, capable of being confidently wrong, which is the
subject of the next two sections.

% ============================================================
\section{Expert 4 --- Adversarial Stress-Testing and Threat Modeling}
\label{sec:redteam}

We evaluate the methodology from an attacker's perspective: not a malicious external
actor, but the more relevant adversary of \emph{systematic self-deception}, in which
the human--AI system converges confidently on a wrong answer. We enumerate the
AI-specific failure modes, note which materialized, and describe the protocols that
did or should guard against them.

\subsection{Confirmation-bias feedback loops}

\textbf{Threat.} An AI assistant optimized for helpfulness can preferentially
produce evidence for whatever hypothesis the human appears to favor, and the human,
seeing agreement, grows more confident, prompting further confirmatory work. The
loop can lock a collaboration onto a false conclusion with mutually reinforcing
certainty.

\textbf{Materialization.} This risk was live throughout. The synthetic-data phase
is the clearest near-miss: an encouraging early result was, briefly, believed.

\textbf{Mitigation used.} The null-validation discipline and the insistence on
controlled experiments repeatedly injected disconfirming pressure. The decisive
correction---re-running on real data---was an act of deliberate self-challenge. The
recommended protocol is to \emph{schedule} disconfirmation: for every headline
result, require an independent computation whose natural outcome would falsify it.

\subsection{Model hallucination and fabricated specifics}

\textbf{Threat.} Language-model assistants can produce fluent, specific, and
\emph{false} content: non-existent citations, incorrect numerical constants,
plausible-but-wrong formulae, or fabricated intermediate values.

\textbf{Materialization.} Two forms occurred. First, bibliographic details were
assembled by the assistant and are known to be unreliable; this is explicitly
flagged in the reproducibility statement and the physics manuscript, with
verification against primary sources listed as a required pre-submission step.
Second, and more insidiously, there was at least one instance of a
fabricated intermediate value entering a hand-assembled table, later caught by
re-derivation from source.

\textbf{Mitigation used.} The end-to-end, self-validating pipeline was the primary
defense: any number that mattered was regenerated from raw data by a script that
also checked its own null, rather than being trusted as an assistant-reported
scalar. The transferable rule is \textbf{no load-bearing hand-carried numbers}: if
a value appears in a conclusion, a script must reproduce it from inputs.

\subsection{Alignment / sycophancy drift}

\textbf{Threat.} Over a long interaction, an assistant may drift toward telling the
human what they want to hear---softening negative results, overstating the
significance of findings, or declining to push back on a flawed request.

\textbf{Materialization.} The pressure existed but was actively countered. Several
results in this project were \emph{deflating} (supernovae are nearly blind; the
transition width is not recoverable; no dark-energy model beat the standard one;
the method does not resolve the Hubble tension). That these were reported plainly,
rather than dressed up, is the relevant evidence that drift was contained.

\textbf{Mitigation used.} Honesty about negative and null results was treated as a
first-class objective, and claims were continuously re-calibrated to the
computations. The protocol we recommend is to make ``report the disappointing
version'' an explicit standing instruction, and to audit late-stage outputs for
tone inflation relative to early-stage ones.

\subsection{Prompt injection and pipeline integrity}

\textbf{Threat.} In workflows where the assistant ingests external content (web
pages, files, tool output), adversarial instructions embedded in that content could
in principle redirect the research or corrupt results.

\textbf{Materialization.} The project ingested literature via search and fetched
real datasets from public repositories. No injection was observed, but the attack
surface was real: a compromised data file or a poisoned web result could have
altered a covariance matrix or a cited value.

\textbf{Mitigation used and recommended.} Data provenance was restricted to known
public repositories, and the covariances were sanity-checked against published
scalar results (e.g.\ reproducing a known published constraint to three significant
figures before use). The general protocol is provenance pinning plus an independent
consistency check on every externally sourced quantity.

\subsection{Linearization and approximation risk}

\textbf{Threat.} An AI will readily deploy a standard tool (here, the Fisher matrix)
without volunteering that its assumptions are violated. Fisher forecasts are linear
approximations, and the features studied here are large and non-linear.

\textbf{Mitigation used.} The Fisher criterion was used only as an
\emph{explanatory} lens for degeneracy directions, with fully non-linear Monte-Carlo
rates retained as the primary results---and this limitation is stated in the
manuscript. The lesson is to demand, for every convenient approximation, an explicit
statement of where it breaks and an independent estimate that does not rely on it.

\subsection{Red-team summary}

\begin{table}[h]
\centering
\caption{AI-specific failure modes, whether they materialized, and the primary
mitigation.}
\label{tab:redteam}
\begin{tabular}{@{}p{0.26\textwidth}p{0.16\textwidth}p{0.5\textwidth}@{}}
\toprule
Failure mode & Materialized? & Primary mitigation \\
\midrule
Confirmation-bias loop & Near-miss & Scheduled disconfirmation; real-data re-run. \\
Hallucinated citations & Yes & Flagged; primary-source verification required pre-submission. \\
Fabricated intermediate value & Yes (caught) & Self-validating pipeline; no hand-carried numbers. \\
Sycophancy / tone drift & Contained & Negative results reported plainly; tone audit. \\
Prompt injection & Not observed & Provenance pinning; consistency checks on external data. \\
Approximation overreach & Yes (bounded) & Non-linear MC retained as primary; Fisher explanatory only. \\
\bottomrule
\end{tabular}
\end{table}

% ============================================================
\section{Expert 3 --- Limitations and Broader Impacts}
\label{sec:limitations}

\subsection{Limitations of the AI-assisted process}

Beyond the failure modes of Section~\ref{sec:redteam}, several structural
limitations deserve emphasis. The workflow's throughput can outpace human
comprehension: it is possible to accumulate correct results faster than any single
person fully internalizes them, creating a comprehension debt that undermines the
researcher's ability to defend or extend the work. The assistant's fluency can also
lend unwarranted authority to provisional claims; text that reads as authoritative
is not thereby correct. Finally, the assistant's knowledge has a training cutoff and
is uneven across subfields, so its literature synthesis is a starting point for human
verification, never a substitute for it.

\subsection{Ethical blind spots}

Two blind spots merit naming. First, \textbf{automation bias}: the ease and polish
of AI-generated output can suppress the healthy skepticism a researcher would apply
to their own hand computation, and this effect is strongest exactly when the output
agrees with prior expectations. Second, \textbf{diffuse accountability}: when an AI
drafts analysis, code, and prose, the locus of scientific responsibility can blur.
We take the position that responsibility remains entirely with the human authors;
the assistant is an instrument, and its errors are the authors' errors to catch.

\subsection{Broader societal impacts}

The positive impact of workflows like this one is a genuine lowering of the barrier
to rigorous computational science, enabling smaller teams and independent
researchers to undertake work previously requiring large collaborations. The
countervailing risk is a potential flood of plausible, fluent, under-verified
manuscripts that strains peer review and pollutes the literature. The mitigating
factor within our control is the same discipline described throughout: results that
are regenerable from raw inputs by validated code are cheap for reviewers to check,
whereas hand-asserted results are not. We suggest that \emph{regenerability} become
an explicit expectation for AI-assisted computational submissions.

\subsection{Reproducibility Statement}
\label{sec:repro}

The scientific project was designed for end-to-end reproducibility. All headline
quantities are regenerated from raw public data by self-validating scripts that
check their own null distributions. We provide the following (placeholders to be
completed at release):

\begin{itemize}[leftmargin=1.5em]
  \item \textbf{Code repository:} \texttt{[REPOSITORY\_URL\_PLACEHOLDER]}
        --- pipeline, model, and analysis scripts, organized by function, with a
        manifest mapping each script to the result it produces.
  \item \textbf{Datasets:} Type Ia supernovae, \texttt{[PANTHEON+\_DOI\_PLACEHOLDER]};
        anisotropic BAO, \texttt{[DESI\_DR2\_DOI\_PLACEHOLDER]}; cosmic-chronometer
        compilation, \texttt{[CC\_COMPILATION\_DOI\_PLACEHOLDER]}. All are public;
        provenance is restricted to the originating repositories.
  \item \textbf{Derived data products:} result arrays and figures,
        \texttt{[ARCHIVE\_DOI\_PLACEHOLDER]}.
  \item \textbf{Environment:} Python
        \texttt{[PYTHON\_VERSION]}; \texttt{numpy [VERSION]},
        \texttt{scipy [VERSION]}, \texttt{matplotlib [VERSION]}. Random seeds for
        all Monte-Carlo runs are fixed and recorded in
        \texttt{[SEED\_MANIFEST\_PLACEHOLDER]}.
  \item \textbf{Validation:} each pipeline script prints its null-distribution
        median and compares it to the analytic $\chi^2_1$ value prior to reporting
        detection rates; reviewers can confirm calibration without rerunning the
        full Monte Carlo.
  \item \textbf{Known verification debt:} bibliographic metadata assembled during
        the project must be checked against primary sources before archival; the
        multi-model CMB treatment uses distance priors rather than a full
        likelihood; one model's high-redshift extension is numerically unreliable
        and only its low-redshift result is validated. These are documented in the
        physics manuscript.
\end{itemize}

% ============================================================
\section{Synthesis and Conclusions}
\label{sec:conclusion}

Read together, the five perspectives converge on a single thesis. The
distinguishing feature of this AI-assisted methodology was not speed for its own
sake but a \emph{relocation of effort}: because execution became nearly free, the
scarce resource became verification, and the quality of the work was determined
almost entirely by how rigorously the human--AI system challenged its own outputs.
Every genuine advance in the project---the grid that fixed the nesting bug, the
real-data reversal that overturned the supernova result, the controlled experiment
that proved the mechanism---was an act of disconfirmation. Every serious risk---the
confirmation loop, the hallucinated specifics, the drift toward comfortable
conclusions---was an act of unchecked confirmation waiting to happen.

The practical recommendation that follows is uncomfortable but simple. In
AI-assisted research, the assistant's fluency is not the product; the verification
scaffolding around it is. A workflow that generates results faster than it can
falsify them is not doing science faster---it is accumulating risk faster. The
methodology documented here worked to the extent that it inverted this: it made
disconfirmation cheap, scheduled, and habitual, and it refused to let any
conclusion rest on a number that a validated script could not reproduce from raw
data. We offer that inversion, rather than any particular result, as the
transferable contribution.

% ============================================================
% Expert 5 --- Bibliography (placeholder BibTeX; verify before use)
\bibliographystyle{plainnat}
\begin{thebibliography}{99}

\bibitem[Chevallier \& Polarski(2001)]{CPL2001}
Chevallier, M. \& Polarski, D. (2001).
\newblock Accelerating universes with scaling dark matter.
\newblock \emph{Int.\ J.\ Mod.\ Phys.\ D}, \textbf{10}, 213.
\newblock [PLACEHOLDER --- verify against primary source.]

\bibitem[Linder(2003)]{Linder2003}
Linder, E.~V. (2003).
\newblock Exploring the expansion history of the universe.
\newblock \emph{Phys.\ Rev.\ Lett.}, \textbf{90}, 091301.
\newblock [PLACEHOLDER --- verify.]

\bibitem[Linder \& Huterer(2005)]{LinderHuterer2005}
Linder, E.~V. \& Huterer, D. (2005).
\newblock How many dark energy parameters?
\newblock \emph{Phys.\ Rev.\ D}, \textbf{72}, 043509.
\newblock [PLACEHOLDER --- verify.]

\bibitem[Moresco et al.(2022)]{Moresco2022}
Moresco, M. et al. (2022).
\newblock Unveiling the universe with cosmic chronometers.
\newblock \emph{Living Reviews in Relativity} [PLACEHOLDER --- verify volume/page].

\bibitem[Scolnic et al.(2022)]{PantheonPlus2022}
Scolnic, D. et al. (2022).
\newblock The Pantheon+ analysis: the full dataset and light-curve release.
\newblock \emph{Astrophys.\ J.} [PLACEHOLDER --- verify].

\bibitem[DESI Collaboration(2025)]{DESIDR2}
DESI Collaboration (2025).
\newblock DESI DR2 results: baryon acoustic oscillations.
\newblock [PLACEHOLDER --- verify arXiv/journal reference.]

\bibitem[Saridakis(2020)]{Saridakis2020}
Saridakis, E.~N. (2020).
\newblock Barrow holographic dark energy.
\newblock \emph{Phys.\ Rev.\ D}, \textbf{102}, 123525.
\newblock [PLACEHOLDER --- verify.]

\bibitem[Chen et al.(2019)]{Chen2019}
Chen, L., Huang, Q.-G. \& Wang, K. (2019).
\newblock Distance priors from Planck final release.
\newblock \emph{JCAP}, \textbf{02}, 028.
\newblock [PLACEHOLDER --- verify.]

\bibitem[Amann et al.(placeholder)]{ReproPlaceholder}
[Author(s)] ([Year]).
\newblock [Foundational reference on reproducibility / research-integrity in
computational science --- to be supplied.]

\end{thebibliography}

\end{document}
